\documentclass[hyperref={pdfpagelabels=false}]{beamer}
% beamer already loads hyperref; do NOT load \usepackage{hyperref} again.

\usepackage{ragged2e} % 文字向兩邊對齊
\renewcommand{\raggedright}{\leftskip=0pt \rightskip=0pt plus 0cm}

\usepackage{listings}
\usepackage{lmodern}
\usepackage{fontspec}
\usepackage{xeCJK}
\setCJKmainfont{Noto Sans CJK TC}
%（可選）若你會用到無襯線中文或等寬中文，建議補上：
\setCJKsansfont{Noto Sans CJK TC}
\setCJKmonofont{Noto Sans Mono CJK TC}

\usepackage[english]{babel}
\usepackage{amsmath,amssymb,amsfonts,amsthm}

% hyperref settings (since beamer already loads hyperref)
\hypersetup{
  colorlinks=true,
  urlcolor=blue,
  citecolor=blue,
  pdfborder={0 0 0}
}

\usepackage{accsupp}% http://ctan.org/pkg/accsupp
\newcommand{\emptyaccsupp}[1]{\BeginAccSupp{ActualText={}}#1\EndAccSupp{}}

\usepackage{tikz}
\usetikzlibrary{arrows,decorations.pathmorphing,backgrounds,fit,positioning,shapes.symbols,chains}

\usepackage{booktabs}
\usepackage{threeparttable}
\usepackage{colortbl} 
\usepackage{xcolor}
\usetikzlibrary{tikzmark}

\newcommand{\indep}{\perp\!\!\!\perp}
\newcommand{\nindep}{\not\!\perp\!\!\!\perp}

\beamertemplatefootpagenumber



\title[]  
%{Introduction to Causal Inference and Big Data}
{Potential Outcomes Framework}

%\subtitle
%{免除部分負擔對幼兒醫療利用與健康的影響:斷點迴歸分析}

%\author 
 %{Tzu-Ting Yang\inst{1} \and Hsing-Wen Han\inst{2} \and Hsien-Ming Lien\inst{3}}

%\institute[short]{\inst{1} Academia Sinica \and \inst{2} Tamkang University \and \inst{3} National Chengchi University}


\author 
{Prof. Tzu-Ting Yang  \\ 楊子霆}

%\institute{Institute of Economics, Academia Sinica  \\ Econ, NTU}
\institute{Institute of Economics, Academia Sinica  \\ 中央研究院經濟研究所}
%\institute{Institute of Economics, Academia Sinica  \\ IMES/IDAS, NCCU}
%\author 
 %{楊子霆 \\ 中研院經濟所}

%\institute 
%{中正大學國際經濟專題研討}




\date
 {\today}


% \usetheme{CambridgeUS}
% %\usecolortheme{seahorse}
% \setbeamertemplate{footline}{%
% 	\hfill\usebeamercolor[fg]{page number in head/foot}%
% 	\usebeamerfont{page number in head/foot}%
% 	\insertframenumber\,/\,\inserttotalframenumber\hspace{1em}\vspace{2pt}%
% }

% \usetheme{Boadilla}
% %\useinnertheme{rectangles}

% %\usetheme{CambridgeUS}
% \usecolortheme{seahorse}
% \setbeamertemplate{footline}{%
% 	\hfill\usebeamercolor[fg]{page number in head/foot}%
% 	\usebeamerfont{page number in head/foot}%
% 	\insertframenumber\,/\,\inserttotalframenumber\hspace{1em}\vspace{2pt}%
% }


\usepackage{beamerthemeshadow}
%\usecolortheme{seahorse}
\setbeamertemplate{footline}{%
	\hfill\usebeamercolor[fg]{page number in head/foot}%
	\usebeamerfont{page number in head/foot}%
	\insertframenumber\,/\,\inserttotalframenumber\hspace{1em}\vspace{2pt}%
}

%\usepackage{beamerthemeshadow}
\beamersetuncovermixins{\opaqueness<1>{25}}{\opaqueness<2->{15}}

%\usecolortheme[named=Blue]{structure}
%\usetheme{CambridgeUS}
%\usepackage{beamerthemeshadow}
%\beamersetuncovermixins{\opaqueness<1>{25}}{\opaqueness<2->{15}}

\begin{document}


\begin{frame}{}
 \titlepage
\end{frame}





\begin{frame}{Causal Effect and Potential Outcomes Framework}
	\begin{itemize} 
		\item Estimating causal effect of treatment is a challenging task
		\vspace{0.2cm} 
		\begin{itemize} 	
		\item Because \textbf{we can NOT observe counterfactual outcomes} if one had chosen different treatments
		\vspace{0.2cm}
			\end{itemize}
		\item In order to obtain causal effect, we need to compare \textbf{observed outcomes} with \textbf{counterfactual outcomes}
		\vspace{0.2cm}
		\item The \textbf{potential outcomes framework} provides a way to think about causal effects in a structured way
		
	\end{itemize}
	
\end{frame}






\title[]  
{Potential Outcomes Framework}


\author 
{}

%\institute{Institute of Economics, Academia Sinica  \\ Econ, NTU}
%\institute{Institute of Economics, Academia Sinica  \\ IMES/IDAS, NCCU}
\institute{}

%\author 
%{楊子霆 \\ 中研院經濟所}

%\institute 
%{中正大學國際經濟專題研討}




\date
{}


\begin{frame}{}
	\titlepage
\end{frame}



\begin{frame}{Potential Outcomes Framework}{Treatment Status}
	\begin{itemize} 
		\item The potential outcomes framework is developed by statistician Donald Rubin
		\vspace{0.2cm}	
			\begin{itemize} 
		\item Rubin, Donald. 1974. \textbf{Estimating Causal Effects of Treatments in Randomized and Nonrandomized Studies}. \textit{Journal of Educational Psychology}, 66(5).
			\end{itemize}
	\end{itemize}
\end{frame}




\begin{frame}{Potential Outcomes Framework}{Treatment Status}
\begin{itemize} 

\item \textcolor{blue}{Treatment}: the intervention/variable whose effect we are interested in
\vspace{0.2cm}		
\item $D_i$: a dummy that indicate whether individual $i$ receive treatment or not
\begin{equation*}
D_i=\left\lbrace
\begin{array}{l}
1\quad \text{if individual $i$ received the treatment} \\
0\quad \text{otherwise.}
\end{array} \right.
\end{equation*}
\item Examples:
\vspace{0.2cm}
\begin{itemize}
\item Attend graduate school or not
\vspace{0.2cm}
\item Have health insurance or not
\vspace{0.2cm}
\item Win a lottery or not
\vspace{0.2cm}
\item Increase corporate tax rate or not
\vspace{0.2cm}
\item Democracy v.s. Dictatorship
\end{itemize}

\end{itemize}
\end{frame}



\begin{frame}{Potential Outcomes Framework}{Treatment Status}
\begin{itemize} 
\item $D_i$ can be a multiple valued (continuous) variable

\begin{equation*}
D_i= s
\end{equation*}
\item Examples:
\vspace{0.2cm}
\begin{itemize}
\item Schooling years
\vspace{0.2cm}
\item Number of children
\vspace{0.2cm}
\item Number of polices
\vspace{0.2cm}
\item Number of advertisements
\vspace{0.2cm}
\item Money supply 
\vspace{0.2cm}
\item Income tax rate
\end{itemize}
\vspace{0.2cm}
\item \textbf{In the following slides, we focus on the case when a treatment $D_i$ is a binary variable}
\end{itemize}
\end{frame}






\begin{frame}{Potential Outcomes Framework}{Potential Outcome}
\begin{itemize} 	
\item A \textcolor{blue}{potential outcome} is the outcome that would be realized according to which treatment an individual received
\vspace{0.2cm}
\item Suppose there are \textbf{two treatments} for each individual: 
\vspace{0.2cm}
\begin{itemize} 	
\item $D_i=1$
\vspace{0.2cm}
\item $D_i=0$ 
\end{itemize}
\vspace{0.2cm}
\item Thus, each individual $i$ has \textbf{two potential outcomes} and one for each value of the treatment	$\mathrm{Y}^{D}_{i}$
\begin{itemize}
\vspace{0.2cm}
\item $\mathrm{Y}^{1}_{i}$: Potential outcome for an individual $i$ if getting treatment
\vspace{0.2cm}
\item $\mathrm{Y}^{0}_{i}$: Potential outcome for an individual  $i$ if not getting treatment
\vspace{0.2cm}
\end{itemize}

\end{itemize}

\end{frame}





\begin{frame}{Potential Outcomes Framework}{Potential Outcome}
\begin{itemize} 	
\item \textbf{Example:}
\vspace{0.2cm}
\begin{itemize} 
\item Annual earnings if attending graduate school
\vspace{0.2cm}
\item Annual earnings if not attending graduate school
\vspace{0.2cm}
\end{itemize}
%\vspace{0.2cm}
%\item Each potential outcome corresponds to the various levels of a treatment

\vspace{0.2cm}
\item Again, potential outcome can be $\mathrm{Y}^{s}_{i}$:
\begin{itemize}
	\vspace{0.2cm}
\item $s$ can be continuous	
	 \vspace{0.2cm}
\item More than two potential outcomes
\vspace{0.2cm}
\end{itemize}
\item \textbf{How many treatments we have, how many potential outcomes will be}

\end{itemize}

\end{frame}







\begin{frame}{Potential Outcomes Framework}{Observed Outcome}

%\textcolor{blue}{Treatment Effect} \\
%\begin{itemize} 
%\item The treatment effect or causal effect of the treatment on the outcome for unit $i$ is the difference between its two potential outcomes:
%\end{itemize}
%	\begin{equation*}
%	\mathrm{Y}^{1}_{i}-\mathrm{Y}^{0}_{i}
%	\end{equation*}	
%\item No one can be treated and untreated at the same time
\begin{itemize} 
\item \textcolor{blue}{Observed outcome}: for each particular individual, we only can observe one potential outcome
\vspace{0.2cm}
\item Observed outcome $\mathrm{Y}_i$ is realized as 

\begin{gather*}
	\mathrm{Y}_i = \mathrm{Y}^{1}_{i}D_i+\mathrm{Y}^{0}_{i}(1-D_i) \\
	\text{or} \\
	\mathrm{Y}_i = \left\lbrace
	\begin{array}{l}
		\mathrm{Y}^{1}_{i}\quad \text{if $D_i=1$} \\
		\mathrm{Y}^{0}_{i}\quad \text{if $D_i=0$}
	\end{array} \right.
\end{gather*}


\item \textbf{Only one potential outcome can be realized}	
\vspace{0.2cm}
\item The unobserved outcome is called the ``\textbf{counterfactual}" outcome
\end{itemize}

\end{frame}








\title[]  
{Casual Effects}


\author 
{}

%\institute{Institute of Economics, Academia Sinica  \\ Econ, NTU}
%\institute{Institute of Economics, Academia Sinica  \\ IMES/IDAS, NCCU}
\institute{}

%\author 
%{楊子霆 \\ 中研院經濟所}

%\institute 
%{中正大學國際經濟專題研討}




\date
{}


\begin{frame}{}
	\titlepage
\end{frame}







\begin{frame}{Casual Effect}
\begin{itemize} 
%\item What are we trying to estimate? 
%\vspace{0.2cm}
\item \textbf{Causal effect:} the comparisons between the
potential outcomes under each treatment
\vspace{0.2cm}  
\begin{itemize} 
\item The differences between \textbf{observed (potential) outcome} and \textbf{counterfactual (potential) outcome} 
\vspace{0.2cm}  
\end{itemize}



\end{itemize}
\end{frame}





\title[]  
{Casual Effect for an Individual}


\author 
{}

%\institute{Institute of Economics, Academia Sinica  \\ Econ, NTU}
%\institute{Institute of Economics, Academia Sinica  \\ IMES/IDAS, NCCU}
\institute{}

%\author 
%{楊子霆 \\ 中研院經濟所}

%\institute 
%{中正大學國際經濟專題研討}




\date
{}


\begin{frame}{}
	\titlepage
\end{frame}



\begin{frame}{Casual Effect for an Individual}
\begin{itemize} 

\item \textcolor{blue} {Individual Treatment Effect (ITE)}:
\begin{equation*}
\tau_i = \mathrm{Y}^{1}_{i} - \mathrm{Y}^{0}_{i}
\end{equation*}		

\begin{itemize} 
\vspace{0.2cm}
%\item Also call  \textcolor{blue} {Individual Causal Effect}
%\vspace{0.2cm}
%\item The difference between an individual $i$'s two potential outcomes
%\vspace{0.2cm}
\item \textbf{Interpretation:}  The difference between an individual $i$'s potential outcome under treatment v.s. without treatment

\vspace{0.2cm}	
\item \textbf{Example:}
\vspace{0.2cm}	
\begin{itemize} 
\item The difference in individual $i$'s earnings if he/she attends graduate school v.s. not attending graduate school
\end{itemize}		

\vspace{0.2cm}
\item We usually cannot identify the ITE


%\item Almost always unidentified without strong assumptions
\end{itemize}


\end{itemize}
\end{frame}







\begin{frame}{Individual Treatment Effect (ITE)}{An Example}
\begin{itemize}
\item Imagine a population with 4 people
\vspace{0.2cm}
\begin{center}
	\begin{threeparttable}
		\centering\footnotesize
		\begin{tabular}{cccccc}
			\toprule
			$i$ & $D_{i}$ & $\mathrm{Y}^{1}_{i}$ & $\mathrm{Y}^{0}_{i}$ & $Y_{i}$ & $\mathrm{Y}^{1}_{i}-\mathrm{Y}^{0}_{i}$  \\
			\midrule
			David & 1 & 3 & \textcolor{red}? & 3 & \textcolor{red}? \\
			Tina & 1 & 2 & \textcolor{red}? & 2 & \textcolor{red}? \\
			Mary & 0 & \textcolor{red}? & 1 & 1 & \textcolor{red}? \\
			Bill & 0 & \textcolor{red}? & 1 & 1 & \textcolor{red}? \\
			\midrule
			$\color{white}E[\mathrm{Y}^{1}_{i}-\mathrm{Y}^{1}_{i}|D_i=1]$ &   &     &  &  &  \\
		\end{tabular}
	\end{threeparttable}
\end{center}


\item We want to evaluate the effect of attending graduate school on the annual earnings
\begin{itemize}
\vspace{0.2cm}
\item $D_{i}$: Attending graduate school $D_{i}=1$, otherwise $D_{i}=0$
\vspace{0.2cm}
\item $\mathrm{Y}^{1}_{i}$: (Potential) annual earnings if individual $i$ attend graduate school
\vspace{0.2cm}
\item $\mathrm{Y}^{0}_{i}$: (Potential) annual earnings if individual $i$ do not attend graduate school
\item $Y_{i}$: Observed annual earnings for individual $i$
\vspace{0.2cm}

\end{itemize}

\end{itemize}
\end{frame}



\begin{frame}{Individual Treatment Effect (ITE)}{An Example}
\begin{itemize}
\item What is Individual causal effect (ITE) of attending graduate school for David?
\vspace{0.2cm}
\begin{itemize}

\item We only observe the annual earnings for David who attended graduate school
\vspace{0.2cm}
\item Only observe $\mathrm{Y}^{1}$
\vspace{0.2cm}
\end{itemize}

\item What is Individual causal effect (ITE) of attending graduate school for Bill?
\vspace{0.2cm}
\begin{itemize}
\item We only observe the annual earnings for Bill who did not attend graduate school
\vspace{0.2cm}
\item Only observe $\mathrm{Y}^{0}$
\end{itemize}

\end{itemize}
\end{frame}




\begin{frame}{Individual Treatment Effect (ITE)}{An Example}
\begin{itemize}
\item Suppose \textbf{we can observe counterfactual outcomes}
\vspace{0.2cm} 
\begin{center}
	\begin{threeparttable}
		\centering\footnotesize
		\begin{tabular}{cccccc}
			\toprule
			$i$ & $D_i$ & $\mathrm{Y}^{1}_{i}$ & $\mathrm{Y}^{0}_{i}$ & $Y_i$ & $\mathrm{Y}^{1}_{i}-\mathrm{Y}^{0}_{i}$  \\
			\midrule
			David & 1 & 3 & \textcolor{red} 2 & 3 & \textcolor{red} 1 \\
			Tina & 1 & 2 & \textcolor{red} 1 & 2 & \textcolor{red} 1 \\
			Mary & 0 & \textcolor{red} 1 & 1 & 1 & \textcolor{red} 0 \\
			Bill & 0 & \textcolor{red} 1 & 1 & 1 & \textcolor{red} 0 \\
			\midrule
			$\color{white}E[Y_1-Y_0|D=1]$ &   &   &   &  &  \\
		\end{tabular}
	\end{threeparttable}
\end{center}

\item The ITE for David: $\alpha_{David}=1$
\vspace{0.2cm}
\item The ITE for Bill: $\alpha_{Bill}=0$


\end{itemize}
\end{frame}






\title[]  
{Causal Effect for General Population}


\author 
{}

%\institute{Institute of Economics, Academia Sinica  \\ Econ, NTU}
%\institute{Institute of Economics, Academia Sinica  \\ IMES/IDAS, NCCU}
\institute{}

%\author 
%{楊子霆 \\ 中研院經濟所}

%\institute 
%{中正大學國際經濟專題研討}




\date
{}


\begin{frame}{}
	\titlepage
\end{frame}




\begin{frame}{Causal Effect for General Population}
%\textcolor{blue}{Problem}\\
\begin{itemize}
\item People might be more interested in the \textbf{causal effect for general population}
\vspace{0.2cm}	
%\item ITE might not represent causal effect for general population
%\vspace{0.2cm}
%\begin{itemize}
%\item We usually cannot rule out that the ITE differs across individuals
%\vspace{0.2cm}
%\begin{itemize}
%\item Heterogeneous treatment effects
%\end{itemize}
%\end{itemize}
\vspace{0.2cm}
\item Understand the treatment effect (causal effect) for general population: 
\vspace{0.2cm}
\begin{itemize}
\item Estimate the \textbf{population average of the individual treatment effects} 
\end{itemize}

\end{itemize}
\end{frame}



\begin{frame}{Review: Expectation}
\begin{itemize} 
\item We usually use $\mathrm{E}[Y_{i}]$ (the expectation of a variable $Y_{i}$) to denote \textbf{population average} of $Y_{i}$ 

\item Suppose we have a population with N individuals

\begin{equation*}
\mathrm{E}[\mathrm{Y}_{i}] = \dfrac{1}{N} \sum^N_{i=1} \mathrm{Y}_{i}
\end{equation*}

\vspace{0.2cm}
%\item If $Y_{i}$ is a variable from a sample survey, $\mathrm{E}[Y_{i}]$ is the average obtained if everyone in the population is surveyed
%\vspace{0.2cm}
%\item For example: average income for whole population in Taiwan

\end{itemize}
\end{frame}






\begin{frame}{Causal Effect for General Population}
\begin{itemize} 

%\item Assume we have a population with $N$ individuals
%\vspace{0.2cm}
\item \textcolor{blue} {Average Treatment Effect (ATE)}:
\begin{equation*}
\alpha_{\text{ATE}} = \mathrm{E}[\tau_i] =  \mathrm{E}[\mathrm{Y}^{1}_{i}-\mathrm{Y}^{0}_{i}] = \dfrac{1}{N} \sum^N_{i=1} [\mathrm{Y}^{1}_{i}-\mathrm{Y}^{0}_{i}]
\end{equation*}

\begin{itemize} 
\vspace{0.2cm}
\item \textbf{Interpretation:} 
\vspace{0.3cm}
\begin{itemize} 
\item Average difference in potential outcomes for the whole population

\end{itemize}
\vspace{0.2cm}
\item \textbf{Example:}
\vspace{0.2cm}
\begin{itemize} 
	\item Average effect of attending graduate school on annual earnings for whole population
	\vspace{0.2cm}
\item Average difference between the earnings of the same individuals if they attend graduate schools v.s. if not attending graduate schools 
\end{itemize}
\vspace{0.2cm}	
\item We'll spend a lot time trying to identify/estimate ATE
\end{itemize}

\end{itemize}
\end{frame}





\begin{frame}{Average Treatment Effect (ATE)}{An Example}
\begin{itemize}
\item Missing data problem also arises when we estimate ATE
\vspace{0.2cm} 
\begin{center}
	\begin{threeparttable}
		\centering\footnotesize
		\begin{tabular}{cccccc}
			\toprule
			$i$ & $D_i$ & $\mathrm{Y}^{1}_{i}$ & $\mathrm{Y}^{0}_{i}$ & $Y_i$ & $\mathrm{Y}^{1}_{i}-\mathrm{Y}^{0}_{i}$  \\
			\midrule
			\textcolor{red}{David} & 1 & 3 & \textcolor{red}? & 3 & \textcolor{red}? \\
			\textcolor{red}{Tina} & 1 & 2 & \textcolor{red}? & 2 & \textcolor{red}? \\
			\textcolor{red}{Mary} & 0 & \textcolor{red}? & 1 & 1 & \textcolor{red}? \\
			\textcolor{red}{Bill} & 0 & \textcolor{red}? & 1 & 1 & \textcolor{red}? \\
			\midrule
			$E[\mathrm{Y}^{1}_{i}]$     &  & \textcolor{red}? &     &  &     \\
			$E[\mathrm{Y}^{0}_{i}]$     &  &     & \textcolor{red}? &  &     \\
			\midrule
			$E[\mathrm{Y}^{1}_{i}-\mathrm{Y}^{0}_{i}]$ &  &     &     &  & \textcolor{red}? \\
			\bottomrule
		\end{tabular}
	\end{threeparttable}
\end{center}

\vspace{0.2cm}
\item What is the effect of attending graduate school on average annual earnings of whole population (ATE)?
\vspace{0.2cm}
\item $\alpha_{\text{ATE}}=E[\mathrm{Y}^{1}_{i}-\mathrm{Y}^{0}_{i}]=?$
\end{itemize}

\end{frame}



\begin{frame}{Average Treatment Effect (ATE)}{An Example}
\begin{itemize}
\item Suppose \textbf{we can observe counterfactual outcomes}
\vspace{0.2cm}
\begin{center}
	\begin{threeparttable}
		\centering\footnotesize
		\begin{tabular}{cccccc}
			\toprule
			$i$ & $D_i$ & $\mathrm{Y}^{1}_{i}$ & $\mathrm{Y}^{0}_{i}$ & $Y_i$ & $\mathrm{Y}^{1}_{i}-\mathrm{Y}^{0}_{i}$  \\
			\midrule
		\textcolor{red}{David} & 1 & 3 & \textcolor{red}{2} & 3 & \textcolor{red}{1} \\
		\textcolor{red}{Tina} & 1 & 2 & \textcolor{red}{1} & 2 & \textcolor{red}{1} \\
		\textcolor{red}{Mary} & 0 & \textcolor{red}{1} & 1 & 1 & \textcolor{red}{0} \\
		\textcolor{red}{Bill} & 0 & \textcolor{red}{1} & 1 & 1 & \textcolor{red}{0} \\
			\midrule
			$E[\mathrm{Y}^{1}_{i}]$     & & 1.75 &     &  &     \\
			$E[\mathrm{Y}^{0}_{i}]$     &     & & 1.25 &  &     \\
			\midrule
			$E[\mathrm{Y}^{1}_{i}-\mathrm{Y}^{0}_{i}]$ &     &     &  &  & 0.5 \\
			\bottomrule
		\end{tabular}
	\end{threeparttable}
\end{center}
\vspace{0.2cm}
\item What is the effect of attending graduate school on average annual earnings of whole population (ATE)?
\vspace{0.2cm}

\item $\alpha_{\text{ATE}}= \dfrac{1+1+0+0}{4} = 0.5$
\end{itemize}

\end{frame}





\title[]  
{Causal Effect for a Specific Sub-population}


\author 
{}

%\institute{Institute of Economics, Academia Sinica  \\ Econ, NTU}
%\institute{Institute of Economics, Academia Sinica  \\ IMES/IDAS, NCCU}
\institute{}

%\author 
%{楊子霆 \\ 中研院經濟所}

%\institute 
%{中正大學國際經濟專題研討}




\date
{}


\begin{frame}{}
	\titlepage
\end{frame}



\begin{frame}{Review: Conditional Expectation}
	\begin{itemize} 
		\item We usually use $\mathrm{E}[Y_{i}|X_{i}=1]$ to denote the average of $Y_{i}$ in the population that has $X_{i}=1$
		
		\item Suppose the population has $N_{1}$ individuals with $X=1$
		
		\begin{equation*}
		\mathrm{E}[Y_{i}|X_{i}=1] = \dfrac{1}{N_{1}} \sum_{i:X=1} \mathrm{Y}_{i}
		\end{equation*}
		
		
		\vspace{0.2cm}
		%\item If $Y_{i}$ is a variable from a sample survey, $\mathrm{E}[Y_{i}|D_{i}=d]$ is the average computed when everyone in the population who has $X_{i}=d$ is surveyed
		%\vspace{0.2cm}
%		\item \textbf{Example:} 
%				\vspace{0.2cm}
%			\begin{itemize} 
%		\item  Average income for those who have master degree in Taiwan 
%		\vspace{0.2cm}
%			\end{itemize}
%		\item The above quantities are unique for specific population
	\end{itemize}
\end{frame}





\begin{frame}{Causal Effect for a Specific Sub-population}
\begin{itemize}
\item \textcolor{blue} {Conditional average treatment effect (CATE)} for a subpopulation:
\begin{equation*}
\alpha_{\text{CATE}} = \mathrm{E}[\tau_i|X_i=f] = \mathrm{E}[\mathrm{Y}^{1}_{i}-\mathrm{Y}^{0}_{i}|X_i=f] = \dfrac{1}{\color{red}N_f} \sum_{\color{red}i:X_i=f}[\mathrm{Y}^{1}_{i}-\mathrm{Y}^{0}_{i}]
\end{equation*}

\begin{itemize}
%\vspace{0.2cm}
\item $N_f$ is the number of units in the subpopulation
\vspace{0.2cm}
\item \textbf{Interpretation:} 
\vspace{0.2cm}
\begin{itemize}
\item Average difference in potential outcomes for the specific subgroup
\vspace{0.2cm}
\end{itemize}
\item \textbf{Example:}
\vspace{0.2cm}
\begin{itemize} 

	\item Average effect of attending graduate school on annual earnings for \textbf{female} ($X_{i}=f$)
	\vspace{0.2cm}
\item  Average difference between the earnings of \textbf{female} if they attend graduate schools v.s. if not attending graduate schools
\end{itemize}


\end{itemize}


\end{itemize}


\end{frame}



\begin{frame}{Causal Effect for Treatment Group}
\begin{itemize}
	\vspace{0.2cm}
  \item \textcolor{blue}{Average treatment effect on the treated (ATT)}:
  \begin{equation*}
    \alpha_{\text{ATT}}
    = \mathrm{E}[\tau_i\mid D_i=1]
    = \mathrm{E}[\mathrm{Y}^{1}_{i}-\mathrm{Y}^{0}_{i}\mid D_i=1]
    = \dfrac{1}{\color{red}N_1}\sum_{\color{red}i:D_i=1}\big[\mathrm{Y}^{1}_{i}-\mathrm{Y}^{0}_{i}\big]
  \end{equation*}

  \item Note that ATT is a special case of CATE
	\vspace{0.2cm}

  \item \textbf{Interpretation:}
  \begin{itemize}
	\vspace{0.2cm}
    \item Average difference in potential outcomes for those who were treated
    	\vspace{0.2cm}

  \end{itemize}

  \item \textbf{Example:}
  	\vspace{0.2cm}

  \begin{itemize}
	\vspace{0.2cm}

    \item Average effect of attending graduate school on annual earnings for those attending graduate school ($D_i=1$)
    	\vspace{0.2cm}

    \item Average difference between the earnings of those attending graduate schools vs.\ earnings if they had not attended graduate schools
  \end{itemize}
\end{itemize}
\end{frame}














\begin{frame}{Average Treatment Effect on Treated (ATT)}
\begin{itemize}
\item Missing data problem also arises when we estimate ATT
\vspace{0.2cm} 
\begin{center}
	\begin{threeparttable}
		\centering\footnotesize
		\begin{tabular}{cccccc}
			\toprule
			$i$ & $D_i$ & $\mathrm{Y}^{1}_{i}$ & $\mathrm{Y}^{0}_{i}$ & $Y_i$ & $\mathrm{Y}^{1}_{i}-\mathrm{Y}^{0}_{i}$  \\
			\midrule
			\textcolor{red}{David} & 1 & 3 & \textcolor{red}? & 3 & \textcolor{red}? \\
			\textcolor{red}{Tina} & 1 & 2 & \textcolor{red}? & 2 & \textcolor{red}? \\
			Mary & 0 & ? & 1 & 1 & ? \\
			Bill & 0 & ? & 1 & 1 & ? \\
			\midrule
			$E[\mathrm{Y}^{1}_{i} | D_i=1]$     &  & 2.5 &     &  &     \\
			$E[\mathrm{Y}^{0}_{i} | D_i=1]$     &     & & \textcolor{red}? &  &     \\
			\midrule
			$E[\mathrm{Y}^{1}_{i} - \mathrm{Y}^{0}_{i} | D_i=1]$ &     &     &  &  & \textcolor{red}? \\
			\bottomrule
		\end{tabular}
	\end{threeparttable}
\end{center}

\vspace{0.2cm}
\item What is the effect of attending graduate school on average annual earnings for those who choose to attend graduate school (ATT)?
\vspace{0.2cm}
\item $\alpha_{\text{ATT}}=E[\mathrm{Y}^{1}_{i} - \mathrm{Y}^{0}_{i} | D_i=1]=?$
\end{itemize}

\end{frame}






\begin{frame}{Average Treatment Effect on Treated (ATT)}
\begin{itemize}
\item Suppose \textbf{we can observe counterfactual outcomes}
\vspace{0.2cm}
\begin{center}
	\begin{threeparttable}
		\centering\footnotesize
		\begin{tabular}{cccccc}
			\toprule
			$i$ & $D_i$ & $\mathrm{Y}^{1}_{i}$ & $\mathrm{Y}^{0}_{i}$ & $Y_i$ & $\mathrm{Y}^{1}_{i}-\mathrm{Y}^{0}_{i}$  \\
			\midrule
			\textcolor{red}{David} & 1 & 3 & \textcolor{red}2 & 3 & \textcolor{red}1 \\
			\textcolor{red}{Tina} & 1 & 2 & \textcolor{red}1 & 2 & \textcolor{red}1 \\
			Mary & 0 & 1 & 1 & 1 & 0 \\
			Bill & 0 & 1 & 1 & 1 & 0 \\
			\midrule
			$E[\mathrm{Y}^{1}_{i}| D_i=1]$     & & 2.5 &     &  &     \\
			$E[\mathrm{Y}^{0}_{i}| D_i=1]$     &     & & \textcolor{red}{1.5} &  &     \\
			\midrule
			$E[\mathrm{Y}^{1}_{i}-\mathrm{Y}^{0}_{i}| D_i=1]$ &     &     &  &  & \textcolor{red}1 \\
			\bottomrule
		\end{tabular}
	\end{threeparttable}
\end{center}

\vspace{0.2cm}
\item What is the effect of attending graduate school on average annual earnings of those who choose to attend graduate school (ATT)?
\vspace{0.2cm}
\item $\alpha_{\text{ATT}}= \dfrac{1+1}{2} = 1$
\end{itemize}

\end{frame}





\begin{frame}{Causal Effect for a Control Group}
\begin{itemize}
  \item \textcolor{blue}{Average treatment effect on the untreated (ATU):}
  \begin{equation*}
    \alpha_{\text{ATU}} = \mathrm{E}[\tau_i|D_i=0]
    = \mathrm{E}[\mathrm{Y}^{1}_{i}-\mathrm{Y}^{0}_{i}|D_i=0]
    = \dfrac{1}{\color{red}N_0} \sum_{\color{red}i:D_i=0}\bigl[\mathrm{Y}^{1}_{i}-\mathrm{Y}^{0}_{i}\bigr]
  \end{equation*}

  \item Note that ATU is a special case of CATE
  \vspace{0.2cm}

  \item \textbf{Interpretation:}
  \begin{itemize}
    \item Average difference in potential outcomes for those who were untreated
  \end{itemize}
  \vspace{0.2cm}

  \item \textbf{Example:}
  \begin{itemize}
    \item Average effect of attending graduate school on annual earnings for those NOT attending graduate school ($D_i=0$)
    \vspace{0.2cm}
    \item Average difference between the earnings of those NOT attending graduate school vs.\ earnings if they had attended graduate school
  \end{itemize}
\end{itemize}
\end{frame}












\begin{frame}{Average Treatment Effect on Untreated (ATU)}
	\begin{itemize}
		\item Missing data problem also arises when we estimate ATU
		\vspace{0.2cm} 
	\begin{center}
		\begin{threeparttable}
			\centering\footnotesize
			\begin{tabular}{cccccc}
				\toprule
				$i$ & $D_i$ & $\mathrm{Y}^{1}_{i}$ & $\mathrm{Y}^{0}_{i}$ & $Y_i$ & $\mathrm{Y}^{1}_{i}-\mathrm{Y}^{0}_{i}$  \\
				\midrule
				David & 1 & 3 & ? & 3 & ? \\
				Tina & 1 & 2 & ? & 2 & ? \\
				\textcolor{red}{Mary} & 0 & \textcolor{red}? & 1 & 1 & \textcolor{red}? \\
				\textcolor{red}{Bill} & 0 & \textcolor{red}? & 1 & 1 & \textcolor{red}? \\
				\midrule
				$E[\mathrm{Y}^{1}_{i} | D_i=0]$     & & \textcolor{red}? &     &  &     \\
				$E[\mathrm{Y}^{0}_{i} | D_i=0]$     & &     & 1 &  &     \\
				\midrule
				$E[\mathrm{Y}^{1}_{i} - \mathrm{Y}^{0}_{i} | D_i=0]$ & &     &     &  & \textcolor{red}? \\
				\bottomrule
			\end{tabular}
		\end{threeparttable}
	\end{center}
	
		\vspace{0.2cm}
		\item What is the effect of attending graduate school on average annual earnings for those who choose NOT to attend graduate school (ATU)?
		\vspace{0.2cm}
		\item $\alpha_{\text{ATU}}=E[\mathrm{Y}^{1}_{i} - \mathrm{Y}^{0}_{i} | D_i=0]=?$
	\end{itemize}
	
\end{frame}






\begin{frame}{Average Treatment Effect on Untreated (ATU)}
	\begin{itemize}
		\item Suppose \textbf{we can observe counterfactual outcomes}
		\vspace{0.2cm}
		\begin{center}
			\begin{threeparttable}
				\centering\footnotesize
				\begin{tabular}{cccccc}
					\toprule
					$i$ & $D_i$ & $\mathrm{Y}^{1}_{i}$ & $\mathrm{Y}^{0}_{i}$ & $Y_i$ & $\mathrm{Y}^{1}_{i}-\mathrm{Y}^{0}_{i}$  \\
					\midrule
					David & 1 & 3 & 2 & 3 & 1 \\
					Tina & 1 & 2 & 1 & 2 & 1 \\
				 \textcolor{red}{Mary} & 0 & \textcolor{red}1 & 1 & 1 & \textcolor{red}0 \\
				 \textcolor{red}{Bill} & 0 & \textcolor{red}1 & 1 & 1 & \textcolor{red}0 \\
					\midrule
					$E[\mathrm{Y}^{1}_{i}| D_i=0]$     & & \textcolor{red}1 &     &  &     \\
					$E[\mathrm{Y}^{0}_{i}| D_i=0]$     & &     & 1 &  &     \\
					\midrule
					$E[\mathrm{Y}^{1}_{i}-\mathrm{Y}^{0}_{i}| D_i=0]$ & &     &     &  & \textcolor{red}0 \\
					\bottomrule
				\end{tabular}
			\end{threeparttable}
		\end{center}
		
		\vspace{0.2cm}
		\item What is the effect of attending graduate school on average annual earnings of those who choose NOT to attend graduate school (ATU)?
		\vspace{0.2cm}
		\item $\alpha_{\text{ATU}}= \dfrac{0+0}{2} = 0$
	\end{itemize}
	
\end{frame}





\begin{frame}{Selection into Treatment}
	\begin{itemize}	
		
		\item In this numerical example, we have $\alpha_{\text{ATT}} > \alpha_{\text{ATE}} > \alpha_{\text{ATU}}$
				\vspace{0.2cm}
		\item This may indicate selection into treatment:
				\vspace{0.2cm}
	\begin{itemize}	
		\item Those who benefit most from the treatment (attending graduate school) are most likely to take it 
			\end{itemize}
	\end{itemize}
	
\end{frame}




\begin{frame}{Summary}
	\begin{itemize}	
		
		\item \textcolor{blue} {Individual Treatment Effect (ITE)}:
		\begin{equation*}
				\alpha_{\text{ITE}} = \mathrm{Y}^{1}_{i} - \mathrm{Y}^{0}_{i}
		\end{equation*}	
		
		
		\item \textcolor{blue} {Average treatment effect (ATE)}:
		\begin{equation*}
			\alpha_{\text{ATE}} =  \mathrm{E}[\mathrm{Y}^{1}_{i}-\mathrm{Y}^{0}_{i}] = \dfrac{1}{\color{red}N} \sum_{i}[\mathrm{Y}^{1}_{i}-\mathrm{Y}^{0}_{i}]
		\end{equation*}
		
		
		\item \textcolor{blue} {Conditional average treatment effect (CATE)}:
		\begin{equation*}
			\alpha_{\text{CATE}} =  \mathrm{E}[\mathrm{Y}^{1}_{i}-\mathrm{Y}^{0}_{i}|X_i=f] = \dfrac{1}{\color{red}N_f} \sum_{\color{red}i:X_i=f}[\mathrm{Y}^{1}_{i}-\mathrm{Y}^{0}_{i}]
		\end{equation*}
		
	\end{itemize}
	
\end{frame}




\begin{frame}{Summary}
	\begin{itemize}	
		\item \textcolor{blue} {Average treatment effect on the treated (ATT)}:
		\begin{equation*}
			\alpha_{\text{ATT}} =  \mathrm{E}[\mathrm{Y}^{1}_{i}-\mathrm{Y}^{0}_{i}|D_i=1] = \dfrac{1}{\color{red}N_1} \sum_{\color{red}i:D_i=1}[\mathrm{Y}^{1}_{i}-\mathrm{Y}^{0}_{i}]
		\end{equation*}
		
		
		\item \textcolor{blue} {Average treatment effect on the untreated (ATU)}:
		\begin{equation*}
			\alpha_{\text{ATU}} =  \mathrm{E}[\mathrm{Y}^{1}_{i}-\mathrm{Y}^{0}_{i}|D_i=0] = \dfrac{1}{\color{red}N_0} \sum_{\color{red}i:D_i=0}[\mathrm{Y}^{1}_{i}-\mathrm{Y}^{0}_{i}]
		\end{equation*}
		
	\end{itemize}
	
\end{frame}




\begin{frame}{Which Causal Effect is most Relevant?}
	\begin{itemize}	
		
		\item There is no correct answer to this question
				\vspace{0.2cm}
		\item ITE is the most specific effect		
					\vspace{0.2cm}
				\begin{itemize}	
					\item It is hard to identified
													\vspace{0.2cm}
				\end{itemize}
		\item ATE is the most general parameter
				\vspace{0.2cm}
			\begin{itemize}	
		\item What if we give a treatment to the average person/firm/unit 
			\end{itemize}
						\vspace{0.2cm}
		ATT is often interesting for policy evaluation
								\vspace{0.2cm}
				\begin{itemize}	
			\item Policy makers might want to know the effect on those who took up the policy
		\end{itemize}
									\vspace{0.2cm}
	\item ATU is sometimes interesting for policy evaluation
			\vspace{0.2cm}
		\begin{itemize}	
			\item We may be concerned about the people who did not take up the policy
						\vspace{0.2cm}
			\item How would they be affected if they took up the policy
		\end{itemize}
		
	\end{itemize}
	
\end{frame}










\title[]  
{Fundamental Problem of Causal Inference}


\author 
{}

%\institute{Institute of Economics, Academia Sinica  \\ Econ, NTU}
%\institute{Institute of Economics, Academia Sinica  \\ IMES/IDAS, NCCU}
\institute{}

%\author 
%{楊子霆 \\ 中研院經濟所}

%\institute 
%{中正大學國際經濟專題研討}




\date
{}


\begin{frame}{}
	\titlepage
\end{frame}















\begin{frame}{Fundamental Problem of Causal Inference}
%\textcolor{blue}{Fundamental Problem of Causal Inference} \\
\begin{itemize}

\item We can never \textbf{directly observe} causal effects 
\vspace{0.2cm}
\begin{itemize}
\item ITE, ATE, CATE, ATT or ATU
\end{itemize}

\vspace{0.2cm}
\item Because we can never observe both potential outcomes $(\mathrm{Y}^{1}_{i},\mathrm{Y}^{0}_{i})$ for any individual
\vspace{0.2cm}
\item For someone receiving the treatment ($D_i = 1$)
\vspace{0.2cm}
\begin{itemize}
\item $\mathrm{Y}^{1}_{i}$ is observed 
\vspace{0.2cm}
\item But $\mathrm{Y}^{0}_{i}$ is the \textbf{unobserved} counterfactual outcome
\vspace{0.2cm}
\begin{itemize}
\item It represents what would have happened to an individual $i$ if assigned to control
\end{itemize}
\end{itemize}
\vspace{0.2cm}
\item We need to compare \textbf{potential outcomes}, but we only have
\textbf{observed outcomes}
\vspace{0.2cm}
\item Causal inference is a set of statistical tools that deal with a \textbf{missing data problem}
%\item We do not know:
%\begin{itemize} 
%\item $E[\mathrm{Y}^{1}_{i}|D_i=0]$: Average level of outcome with treatment for those who are not treated
%\item $E[\mathrm{Y}^{0}_{i}|D_i=1]$: Average level of outcome without treatment for those who are treated
%\end{itemize}

%\item Because no one can be treated and untreated at the same time


\end{itemize}

\end{frame}






\title[]  
{Potential Outcome Framework in Economics}


\author 
{}

%\institute{Institute of Economics, Academia Sinica  \\ Econ, NTU}
%\institute{Institute of Economics, Academia Sinica  \\ IMES/IDAS, NCCU}
\institute{}

%\author 
%{楊子霆 \\ 中研院經濟所}

%\institute 
%{中正大學國際經濟專題研討}




\date
{}


\begin{frame}{}
	\titlepage
\end{frame}









\begin{frame}{Potential Outcome Framework in Economics}{Demand and Supply}
	
	
	\begin{itemize} 
		\item The concepts of potential and observed outcomes are deeply ingrained in economics
		\vspace{0.2cm}
		\begin{itemize} 
			\item A demand function represents the potential quantity demanded as a function of price 
			\vspace{0.2cm}
			\item Only the quantity under equilibrium price is realized 
			\vspace{0.2cm}
			\item Other quantities along demand curve are counterfactual			
			
			
			
			
		\end{itemize}
	\end{itemize}
	
\end{frame}




\begin{frame}{Potential Outcome Framework in Economics}{Demand and Supply}
	\begin{center}
		\begin{tikzpicture}[scale=0.7]
		\draw[thick,->] (0,0) --node[very near end,above,sloped]{Price}(0,10);
		\draw[thick,->] (0,0) --node[very near end,below,sloped]{Quantity}(10,0);
		\draw[draw=blue] (0,1) edge [bend right=20] node[very near end,above,sloped]{Supply}(9,9);
		\draw[draw=red] (0,9) edge [bend right=20] node[very near start,above,sloped]{Demand}(9,1);
		\node [below left] at (0,0) {$0$};
		
		\node [below] at (2.5,0) {$Q_1$};
		\node [left] at (0,5.3) {$P_1$};
		\draw[dashed](0,5.3)--(2.5,5.3)--(2.5,0);
		\node [circle , fill, inner sep=1.5pt] at (2.5,5.3) {};
		
		\node [below] at (5.9,0) {$Q_2$};
		\node [left] at (0,2.5) {$P_2$};
		\draw[dashed](0,2.5)--(5.9,2.5)--(5.9,0);
		\node [circle , fill, inner sep=1.5pt] at (5.9,2.5) {};
		\end{tikzpicture}
	\end{center}
\end{frame}







\begin{frame}{Potential Outcome Framework in Economics}{Demand and Supply}
	\begin{center}
		\begin{tikzpicture}[scale=0.7]
		\draw[thick,->] (0,0) --node[very near end,above,sloped]{Price}(0,10);
		\draw[thick,->] (0,0) --node[very near end,below,sloped]{Quantity}(10,0);
		\node [below left] at (0,0) {$0$};
		\node [below] at (4.5,0) {$Q^*$};
		\node [left] at (0,3.4) {$P^*$};
		\draw[draw=blue] (0,1) edge [bend right=20] node[very near end,above,sloped]{Supply}(9,9);
		\draw[draw=red] (0,9) edge [bend right=20] node[very near start,above,sloped]{Demand}(9,1);
		\draw[dashed](0,3.4)--(4.5,3.4)--(4.5,0);
		\tiny{\node [circle , fill, text width=0.5em] at (4.5,3.4) {};}
		\end{tikzpicture}
	\end{center}
\end{frame}




\title[]  
{Stable Unit Treatment Value Assumption}


\author 
{}

%\institute{Institute of Economics, Academia Sinica  \\ Econ, NTU}
%\institute{Institute of Economics, Academia Sinica  \\ IMES/IDAS, NCCU}
\institute{}

%\author 
%{楊子霆 \\ 中研院經濟所}

%\institute 
%{中正大學國際經濟專題研討}




\date
{}


\begin{frame}{}
	\titlepage
\end{frame}







\begin{frame}{Stable Unit Treatment Value Assumption (SUTVA)}
	\begin{block}{Assumption}
		Observed outcomes are realized as 
		\begin{equation*}
		\mathrm{Y}_i=\mathrm{Y}^{1}_{i}D_i+\mathrm{Y}^{0}_{i}(1-D_i)
		\end{equation*}
	\end{block}
	\begin{itemize} 
		\item Implies that observed outcomes for an individual $i$ are \textbf{unaffected} by the treatment status of other individual $j$
		\vspace{0.2cm}
		
		\item Individual $i$'s observed outcomes are only affected by his/her own treatment 
		\vspace{0.2cm}
		\item Rules out possible treatment effect from other individuals (spillover effect/externality)
		
	\end{itemize}
\end{frame}



\begin{frame}{Stable Unit Treatment Value Assumption (SUTVA)}
	
	\begin{itemize} 
		\item  Could write out potential outcomes in a more
		extensive fashion, taking into account both
		one's own treatment status and the treatment
		status of others
		
		\begin{equation*}
		\left\lbrace
		\begin{array}{l}
		\mathrm{Y}^{11}_{i}\quad \text{if $D_i=1$ and $D_j=1$} \\
		\mathrm{Y}^{10}_{i}\quad \text{if $D_i=1$ and $D_j=0$} \\
		\mathrm{Y}^{01}_{i}\quad \text{if $D_i=0$ and $D_j=1$} \\
		\mathrm{Y}^{00}_{i}\quad \text{if $D_i=0$ and $D_j=0$}
		\end{array} \right.
		\end{equation*}
		
		%\begin{equation*}
		%	\mathrm{Y}_i=\mathrm{Y}^{1}_{i}D_i+\mathrm{Y}^{0}_{i}(1-D_i) + \mathrm{Y}^{2}_{i}D_{j} + \mathrm{Y}^{3}_{i}(1-D_{j})
		%	\end{equation*}
		
		\item \textbf{Example:} 
				\vspace{0.2cm}
			\begin{itemize} 
		\item Your health status depends on whether you smoke and your father/mother smoke
			\end{itemize}
		%\item E.g., housemates, friends, relatives, neighbors,
		%competitors…
	\end{itemize}
\end{frame}










\begin{frame}{Stable Unit Treatment Value Assumption (SUTVA)}{Examples for Spillover Effect}
	
		
		\begin{itemize} 
			\vspace{0.2cm}
			\item \textbf{Contagion}: 
						\vspace{0.2cm}
					\begin{itemize} 
			\item The effect of being vaccinated on one's probability of contracting a disease depends on whether others have been vaccinated
					\end{itemize}		
			\vspace{0.2cm}
			\item \textbf{Displacement}: 
			\vspace{0.2cm}
			\begin{itemize} 
			\item Police interventions designed to suppress crime in one location may displace criminal activity to nearby locations.	
								\end{itemize}				
			\vspace{0.2cm} 
			\item \textbf{Communication}:
				\vspace{0.2cm}
			\begin{itemize}  
			\item Interventions that convey information about commercial products, entertainment, or political causes may spread from individuals who receive the treatment to others who are nominally untreated
			\end{itemize}						
		\end{itemize}		
		
		

\end{frame}




\begin{frame}{Stable Unit Treatment Value Assumption (SUTVA)}
	\begin{itemize} 
		
		\item SUTVA may be problematic, so we should choose the units of analysis to minimize interference across units.
		\vspace{0.2cm}
		\item Recent literatures on causal inference are trying to deal with this assumption
		
	\end{itemize}
\end{frame}





\begin{frame}{Suggested Readings}
\begin{itemize} 
\item Chapter 1, Mastering Metrics: The Path from Cause to Effect
\vspace{0.2cm}
\item Chapter 2, Mostly Harmless Econometrics
\vspace{0.2cm}
\item Chapter 4, Causal Inference: The Mixtape
\end{itemize}
\end{frame}





\end{document} 